% sections/technical_notes.tex
% Seção 4: Informações para Integração

\section{Informações para Integração}

Esta seção apresenta a arquitetura do IntelliStock e orientações para
integração por software houses.

\subsection{Arquitetura do Produto}

O IntelliStock é composto por módulos integrados:

\begin{center}
\begin{tabular}{l p{8cm}}
\toprule
\textbf{Módulo} & \textbf{Função} \\
\midrule
Motor de Inteligência & Análise das 4 esferas e geração de alertas \\
API de Integração     & Interface REST para comunicação \\
Persistência Local    & Armazenamento de dados operacionais \\
Dashboard             & Interface web para visualização e gestão \\
\bottomrule
\end{tabular}
\end{center}

\textbf{Fluxo de dados:}

\begin{verbatim}
    [ERP/PDV]  -->  [API de Ingestão]  -->  [Motor de Inteligência]
                                                    |
                                                    v
                                              [Alertas]
                                                    |
                                                    v
                                              [Dashboard]
\end{verbatim}

\subsection{API de Integração}

A API REST permite integração com sistemas externos.

\textbf{Endpoints de Ingestão:}

\begin{itemize}[leftmargin=2em]
  \item \texttt{POST /ingest/items} - Cadastro de itens
  \item \texttt{POST /ingest/stock\_snapshot} - Posição de estoque
  \item \texttt{POST /ingest/sales\_batch} - Vendas consolidadas
  \item \texttt{POST /ingest/dishes} - Cadastro de pratos
  \item \texttt{POST /ingest/recipes} - Receitas (BOM)
  \item \texttt{POST /ingest/lots} - Lotes com validade
\end{itemize}

\textbf{Endpoints de Consulta:}

\begin{itemize}[leftmargin=2em]
  \item \texttt{GET /alerts} - Lista de alertas ativos
  \item \texttt{GET /dashboard/overview} - Resumo do estado atual
  \item \texttt{GET /health} - Status do serviço
  \item \texttt{GET /metrics} - Métricas de análise
\end{itemize}

\textbf{Endpoints de Gestão:}

\begin{itemize}[leftmargin=2em]
  \item \texttt{GET /stock/items} - Lista de itens
  \item \texttt{GET /stock/items/\{id\}} - Detalhes de item
  \item \texttt{PUT /stock/items/\{id\}/count} - Atualização de contagem
  \item \texttt{POST /stock/items/\{id\}/lots} - Adicionar lote
\end{itemize}

\textbf{Documentação interativa:} Disponível em \texttt{/docs} (Swagger UI).

\subsection{Formato de Dados}

Exemplos de payload para integração:

\textbf{Cadastro de Item:}

\begin{verbatim}
{
  "external_id": "SKU001",
  "name": "Filé Mignon",
  "type": "ingredient",
  "unit": "kg",
  "item_class": "A",
  "lead_time_days": 2,
  "shelf_life_days": 5
}
\end{verbatim}

\textbf{Registro de Venda:}

\begin{verbatim}
{
  "sales": [
    {
      "dish_code": "PRT001",
      "quantity": 3,
      "timestamp": "2026-01-03T12:30:00"
    }
  ]
}
\end{verbatim}

\textbf{Registro de Lote:}

\begin{verbatim}
{
  "item_id": "SKU001",
  "lot_id": "LOTE-2026-001",
  "quantity": 10.0,
  "expires_at": "2026-01-10"
}
\end{verbatim}

\subsection{Gestão FEFO}

A Esfera 4 implementa gestão de validade por lote com diferentes níveis
de confiabilidade:

\begin{itemize}[leftmargin=2em]
  \item \textbf{Alta:} Dados de lote informados via API
  \item \textbf{Média:} Validade informada no cadastro de estoque
  \item \textbf{Baixa:} Estimativa baseada em shelf life do item
\end{itemize}

\subsection{Requisitos de Ambiente}

\begin{center}
\begin{tabular}{l l}
\toprule
\textbf{Requisito} & \textbf{Especificação} \\
\midrule
Memória RAM      & 512 MB (mínimo) \\
Disco            & 100 MB + dados \\
Sistema          & Windows, Linux, macOS \\
Rede             & Localhost ou rede interna \\
\bottomrule
\end{tabular}
\end{center}

\subsection{Escopo Atual do Piloto}

O piloto contempla as seguintes funcionalidades:

\begin{itemize}[leftmargin=2em]
  \item Análise completa das 4 esferas
  \item Dashboard operacional
  \item API de integração documentada
  \item Gestão de lotes com FEFO
  \item Configuração de parâmetros por item
\end{itemize}

\textbf{Foco atual:}

\begin{itemize}[leftmargin=2em]
  \item Operação single-unit (uma unidade operacional)
  \item Lead time configurado por item
  \item Análise baseada em quantidade (sem custos financeiros)
  \item Execução em rede local/interna
\end{itemize}

\subsection{Roteiro de Integração}

Para integrar o IntelliStock com um ERP/PDV existente:

\begin{enumerate}[leftmargin=2em]
  \item Mapear itens do ERP para formato de ingestão
  \item Criar rotina de sincronização de estoque (snapshot)
  \item Criar rotina de importação de vendas
  \item Configurar lead time e shelf life por item
  \item Configurar receitas (BOM) para pratos com produção
  \item Testar com dados reais em ambiente isolado
  \item Validar alertas antes de uso em produção
\end{enumerate}

%----------------------------------------------------------------------
\newpage
\section{Planejamento e Próximas Etapas}

Esta seção apresenta o roadmap de evolução do IntelliStock pós-piloto.

\begin{infobox}[Visão de Produto]
Os itens abaixo fazem parte do planejamento estratégico para as próximas
versões do IntelliStock. Não representam funcionalidades entregues no piloto.
\end{infobox}

\subsection{Módulo de Integração (Connector)}

Desenvolvimento de conectores dedicados por software house, permitindo
integração nativa com ERPs e PDVs específicos do mercado.

\textbf{Objetivo:} Simplificar a integração e reduzir esforço de
implantação para parceiros.

\subsection{Empacotamento como Executável}

Distribuição do motor de inteligência como binário executável (.exe),
eliminando dependências de ambiente de desenvolvimento.

\textbf{Objetivo:} Facilitar instalação e atualização em ambientes
de produção.

\subsection{Modelo de Licenciamento}

Estruturação de modelo de licenciamento local para distribuição
comercial através de parceiros.

\textbf{Objetivo:} Viabilizar comercialização e suporte através de
software houses.

\subsection{Endurecimento de Segurança}

Implementação de autenticação, controle de acesso e criptografia
para operação em ambientes de produção.

\textbf{Objetivo:} Atender requisitos de segurança corporativos.

\subsection{Interface Desktop}

Empacotamento do dashboard como aplicativo desktop standalone,
com experiência nativa de usuário.

\textbf{Objetivo:} Oferecer experiência profissional independente
de navegador.

\subsection{Expansão Multi-Unidade}

Suporte a múltiplas unidades operacionais com estoques independentes
e visão consolidada.

\textbf{Objetivo:} Atender redes e operações com mais de uma loja.

\vspace{1em}

\begin{infobox}[Próximos Passos]
As evoluções listadas serão priorizadas com base no feedback do piloto
e nas necessidades identificadas junto a parceiros de integração.
\end{infobox}
